\chapter{The Authority of the Proof Checker}
\label{ch:authority}

The reliability of a proof checker is a technical achievement, and the argument of this monograph depends on it. The question of this chapter is the extent of its authority. A proof checker is an authority on one question, whether a given term has a given type in a given environment, and it is not an authority on the questions that surround that one. The chapter states what the question is, lists what the answer depends on, classifies the ways in which the claims made on the checker's behalf can exceed its jurisdiction, and argues that the checker can be entirely reliable while those claims are not warranted.

\section{What the kernel decides}

The kernel of a system such as Lean checks that a term inhabits a type under an accepted logical environment. In the notation of Chapter~\ref{ch:proofnot} and Appendix~\ref{app:formal}, it decides the judgment $K(V) \iff \Gamma \vdash \pi : F$, where $\Gamma$ is the environment, $F$ the proposition, and $\pi$ the term. The judgment is a function of those three arguments. It does not mention the natural-language specification $N$, the purported argument $P$, or the interpretation $\mathcal{I}$, and Theorem~\ref{thm:rv001} exploits exactly this fact. It does not adjudicate the historical or linguistic intentions behind the type.

A proof assistant is organized so that the kernel is a small trusted component and the rest is not trusted. Tactics, automation, and the elaboration of source syntax into terms are untrusted in the sense that their outputs are rechecked by the kernel, and an error in a tactic can cause a proof to fail and cannot cause a false term to be accepted unless it exploits a defect of the kernel. This is the familiar de Bruijn criterion, and it explains why a proof that the kernel accepts can be trusted as a proof of its type. It also explains the limits of the trust, because everything that precedes the kernel's input, including what the statement says, belongs to the untrusted part and is never checked by the kernel against anything outside the system.

\section{What the answer depends on}

The kernel's judgment is relative to the environment $\Gamma$, and the environment is a substantial object. It contains the axioms in force, the definitions that give meaning to the symbols of the statement, the imported libraries with their theorems and conventions, any constants declared opaque or without a defining equation, and any assumptions the development has introduced. Each of these is a place where a claim made about the verification can outrun what was checked.

\emph{Axioms} fix the logical strength under which acceptance holds. A result accepted in an environment with additional axioms is a result relative to them, and a claim that a theorem was proved in the base system is false if an extra axiom was used. For this reason, the source's criterion for success in translating a proof into Lean includes that Lean accepts the result without a placeholder for a missing proof and without additional axioms. \emph{Definitions} determine what the statement says. A definition may be total by convention in a way that departs from the informal notion, as in the library infimum that returns zero on the empty set, which is the source of the silent default of Chapter~\ref{ch:hidden}. The kernel accepts the definition because it is a definition, and it does not check that it matches the informal object. \emph{Imports} bring in the conventions of a library, and a result that uses an imported lemma is a result relative to the library's statement of the lemma, not to its informal description. \emph{Opaque constants and assumptions} enter the environment without a definition and so contribute to what the final theorem is conditional on, whether or not the theorem's statement displays the dependence.

A practical corollary is that an environment should be reported with a result, since acceptance is acceptance in an environment. The ledger of Chapter~\ref{ch:system} stores it for that reason.

\section{Three kinds of defect and a fourth kind of excess}

The defects that can affect a verified result are of different kinds and the chapter keeps them apart. A \emph{kernel-level defect} is a flaw in the trusted component that allows a false term to be accepted. It is the kind of defect that the engineering of proof assistants aims to exclude, and it is rare. A \emph{specification error} is a discrepancy between the formal statement and the intended statement, and it is not a defect of the kernel, since the kernel correctly checked the statement it was given. It includes the weakening, the lateral substitution, and the silent default documented in earlier chapters. An \emph{environment defect} is a flaw in a definition, an import, or an assumption on which the statement depends, so that the statement is not what it appears to be even though it is stated as intended. These three are defects in the object. The fourth category is an excess in the claim: an \emph{inappropriate interpretation of a verified result}, in which a correct verification of a correct statement is read as establishing more than it does, for example as confirming a particular published argument.

The distinction between the first and the others is the most important for institutional purposes. Criticism of a verification effort is sometimes taken to be criticism of the kernel and is answered with the correct observation that the kernel is reliable. The documented discrepancies of Part~IV are not kernel-level. They are specification-level, and they are untouched by the kernel's reliability. The monograph's statement is that a verifier can be entirely reliable while the claims made on its behalf exceed its jurisdiction, and the examples show that the two conditions are independent.

\section{Jurisdiction}

The word jurisdiction is used in the following sense. The jurisdiction of a verifier is the set of questions on which its output constitutes a certificate in the sense of Appendix~\ref{app:formal}: a statement accepted by an identified authority under an identified interpretation with an identified set of dependencies. For the kernel the authority is the trusted checker, the interpretation is the type theory, and the dependencies are the environment. The certificate is of the formal obligation $O_{\mathrm{formal}}$. It is not a certificate of $O_{\mathrm{statement}}$ or $O_{\mathrm{argument}}$, whose certificates, if any, come from other authorities under other interpretations, such as a human reader's judgment that a definition matches the informal notion.

An assertion of the form that a theorem has been verified therefore has a jurisdictional reading, namely that $O_{\mathrm{formal}}$ is discharged. Other readings are claims about obligations on which the verifier was never consulted. The practice recommended throughout this monograph is that each such claim be attached to the obligation it concerns and to the authority behind it. Where no authority has ruled, the obligation is open, and the honest report says so. The first result of Part~V, Theorem~\ref{thm:rv001}, makes the statement that kernel acceptance does not entail statement correspondence precise, and it is stated here in outline: since the kernel's judgment is invariant under changes to the originating specification, no function of it can determine whether the formal statement corresponds to that specification.
